\section{Discussion}
\label{sec:discussion}

\textbf{Why lightweight routing helps.} The slow or incomplete regimes of the fixed strategies lie at opposite ends of
the selectivity axis: exact search is slow when many rows qualify, iterative HNSW when few do, and single-pass HNSW is
incomplete below about 25\%. A router therefore only needs to know on which side of the crossover a query falls, and
PostgreSQL's histograms answer that with errors of a few percent near the crossover. This is far smaller than the
spacing of the selectivity levels. The router thus closes almost all of the gap to the hindsight latency floor at a
cost of microseconds, which answers RQ2 positively for this workload. The benefit lies in the tail and in robustness
across selectivities, not at individual levels, where P1 is at best on par with the arm it selects.

\textbf{Why the default planner is fast but incomplete.} Choosing HNSW from 2\% upwards is the cheapest plan by
PostgreSQL's cost model, which assumes the scan returns every qualifying row. With a post-filtered single-pass scan
that assumption fails for 27\% of the workload. Enabling iterative scanning would restore completeness, but iterative
scanning changes execution, not costing, so the planner would still pick HNSW at 2\%, where iterative HNSW is about
twice as slow as exact search (7.29 vs.\ 3.46\,ms). This supports~\cite{amanbayev2026filtered} and suggests the
planner would need an estimate of result completeness, or of iterative-scan cost, to choose well.

\textbf{Strengths and weaknesses of a single threshold.} One threshold on the estimated row count needs no model and
no engine change, and it works when the two cost curves are monotone in selectivity and cross once, as they do here.
Its weakness is the crossover itself. At 5\% the faster arm varies between query vectors with the same row count, and
the mean-latency objective accepted a tail loss over a flat interval of thresholds. The threshold is also specific to
the hardware, the HNSW parameters and the attribute distribution. This answers RQ3: P1 loses measurably only near the
crossover (in the tail) and, by a fraction of a millisecond, at the broadest levels where single-pass HNSW is already
complete; nowhere does its overhead outweigh its savings in aggregate. The hypothesis that exact search wins for few qualifying rows and HNSW for many is supported,
with a crossover at 2{,}000--5{,}000 of 100{,}000 rows.

\subsection{Limitations}
The results come from one dataset (100{,}000 LAION/CLIP vectors), one numeric attribute and closed-range predicates.
The attribute is a CLIP score that may correlate with the embedding geometry, so other attributes or categorical
filters may shift the crossover. All measurements use one machine, with PostgreSQL in Docker on a Windows desktop,
which adds latency and noise; absolute numbers will differ on other hardware. The benchmark is warm-cache and
single-client: concurrent workloads, where B1's parallel scans compete for CPU, and cold-cache execution, where HNSW's
random accesses become costlier, were not evaluated. The table is static, so statistics staleness was not studied.
The HNSW parameters were fixed in advance, and the threshold is valid only for this configuration. The held-out split
has 90 query vectors, and per-level differences of a few tens of microseconds are within run-to-run noise. The
improvements should
not be generalized to other data sizes, filter distributions, hardware or workloads without re-calibration.
